% Source problem: 2013ap_exam.pdf, PDF pp. 8--9.
\documentclass[11pt,a4paper]{article}
\usepackage[margin=23mm]{geometry}
\usepackage{fontspec}
\usepackage{xeCJK}
\setCJKmainfont{HaranoAjiMincho-Regular.otf}
\setCJKsansfont{HaranoAjiGothic-Medium.otf}
\usepackage{amsmath,amssymb}
\setlength{\parindent}{1em}
\setlength{\parskip}{0.35em}
\title{物理工学専攻 院試解答集\\2013年 物理学II 第2問}
\author{政岡凜太郎}
\date{}
\begin{document}
\maketitle

\section*{[1] 速度空間でのミクロカノニカル分布}
$N=2$ の全運動エネルギーは
\begin{equation}
 E=\frac{m}{2}(v_{x,1}^2+v_{y,1}^2+v_{x,2}^2+v_{y,2}^2).
\end{equation}
したがって4成分を座標とする速度空間では、許される点は原点中心、半径
\begin{equation}
 \boxed{R=\sqrt{\frac{2E}{m}}}
\end{equation}
の3次元球面上にある。ミクロカノニカル分布では、等しい微小エネルギー幅の殻に含まれる位相空間体積を等重みに数える。
速度球面の面積要素を $dS$、半径方向の厚さを $dR$ とすると、殻の体積は $dS\,dR$ である。
$E=mR^2/2$ を微分すれば
\begin{equation}
  dE=mR\,dR,\qquad dS\,dR=\frac{dS\,dE}{mR}.
\end{equation}
同じ球面上では $mR$ が一定なので、固定したエネルギー幅に対する重みは $dS$ に比例する。
位置積分は速度方向に依存しない共通因子となり、ミクロカノニカル分布はこの球面に一様である。

\section*{[2] $d$ 次元球面の表面積}
半径 $r$ の球面を、最初の $d-1$ 座標 $q_1,\ldots,q_{d-1}$ で表すと
\begin{equation}
 q_d=\pm\sqrt{r^2-q_1^2-\cdots-q_{d-1}^2}.
\end{equation}
各半球の面積要素はグラフ面積の公式から
\begin{equation}
 dS=\sqrt{1+\sum_{j=1}^{d-1}\left(\frac{\partial q_d}{\partial q_j}\right)^2}
       \,dq_1\cdots dq_{d-1}
    =\frac{r}{\sqrt{r^2-\sum_{j=1}^{d-1}q_j^2}}\,dq_1\cdots dq_{d-1}.
\end{equation}
上下二つの半球を足し、各変数が残りの半径以内を動く積分にすると、問題文の式が得られる。たとえば
\begin{align}
 S_2(r)&=\int_{-r}^{r}dq_1\,\frac{2r}{\sqrt{r^2-q_1^2}},\\
 S_3(r)&=\int_{-r}^{r}dq_1
 \int_{-\sqrt{r^2-q_1^2}}^{\sqrt{r^2-q_1^2}}dq_2\,
 \frac{2r}{\sqrt{r^2-q_1^2-q_2^2}},\\
 S_4(r)&=\int_{-r}^{r}dq_1
 \int_{-\sqrt{r^2-q_1^2}}^{\sqrt{r^2-q_1^2}}dq_2
 \int_{-\sqrt{r^2-q_1^2-q_2^2}}^{\sqrt{r^2-q_1^2-q_2^2}}dq_3\,
 \frac{2r}{\sqrt{r^2-q_1^2-q_2^2-q_3^2}}.
\end{align}
同じ構成を $d-1$ 重積分へ拡張できる。

\section*{[3] $N=2$ の粒子1の速さ}
速度ベクトル $\boldsymbol v_1$ が平面内で速さ $v_1$ と角度 $\theta$ を持つとする。エネルギー一定面上で粒子2の速さは
\begin{equation}
 s=|\boldsymbol v_2|=\sqrt{R^2-v_1^2}.
\end{equation}
各 $\boldsymbol v_1$ に対し、残る2成分 $(v_{x,2},v_{y,2})$ は半径 $s$ の円周上を動く。エネルギーのデルタ関数を用いてその寄与を書くと
\begin{align}
 P(v_1)\,dv_1
 &\propto \underbrace{2\pi v_1\,dv_1}_{\boldsymbol v_1\text{ の円環面積}}
 \int_{\mathbb R^2}d^2v_2\,
 \delta\!\left(E-\frac m2(v_1^2+v_2^2)\right)\\
 &\propto 2\pi v_1\,dv_1\,\frac{S_2(s)}{m s}.
\end{align}
$0<v_1<R$ で粒子2の速さを積分変数 $u$ とすると、デルタ関数の引数が0となるのは $u=s$ で、その微分の絶対値は $ms$ である。
したがって上の殻積分は
\begin{equation}
  \int_0^\infty 2\pi u\,du\,
  \delta\!\left(E-\frac m2(v_1^2+u^2)\right)
  =\frac{2\pi s}{ms}=\frac{S_2(s)}{ms}=\frac{2\pi}{m}.
\end{equation}
ここで $S_2(s)=2\pi s$ は円周長であり、この殻積分は $v_1$ に依存しない。したがって速さの確率密度は
\begin{equation}
 \boxed{P(v_1)\propto v_1\quad(0<v_1<R),\qquad
 P(v_1)=0\quad(v_1>R).}
\end{equation}
半径 $v_1$ の円環面積に比例するため、平面速度ベクトルの密度と速さの密度は異なる。

\section*{[4] 一般の $N$ に対する速さ分布}
全エネルギー $E=N\epsilon$ のとき、$2N$ 次元速度球面の半径は
\begin{equation}
 R=\sqrt{\frac{2N\epsilon}{m}}.
\end{equation}
粒子1の二成分を固定すると、残る $2N-2$ 成分の速度空間は半径
$s=\sqrt{R^2-v_1^2}$ の球面上にある。前問と同じエネルギー殻積分から
\begin{equation}
 P(v_1)\,dv_1\propto 2\pi v_1\,dv_1\,
 \frac{S_{2N-2}(s)}{m s}.
\end{equation}
$d$ 次元球面の表面積は半径の $d-1$ 乗に比例するので
$S_{2N-2}(s)\propto s^{2N-3}$。従って
\begin{equation}
 \boxed{P(v_1)\propto v_1
 \left(\frac{2N\epsilon}{m}-v_1^2\right)^{N-2},\qquad
 0<v_1<\sqrt{\frac{2N\epsilon}{m}}.}
\end{equation}
それより大きい速さでは $P(v_1)=0$ である。

\section*{[5] Maxwell速さ分布への極限}
$v_1$ を固定して $N\to\infty$ とする。$v_1$ に依らない係数を比例定数へまとめると
\begin{align}
 P(v_1)&\propto v_1\left(1-\frac{mv_1^2}{2N\epsilon}\right)^{N-2}\\
 &\longrightarrow v_1\exp\!\left(-\frac{mv_1^2}{2\epsilon}\right).
\end{align}
これは二次元でのMaxwell速さ分布であり、先頭の $v_1$ は二次元の円環面積要素に由来する。

\section*{[6] 温度}
温度 $T$ の二次元Maxwell速さ分布は
$P(v_1)\propto v_1\exp[-mv_1^2/(2k_BT)]$ である。指数を[5]と比べると
\begin{equation}
 \boxed{T=\frac{\epsilon}{k_B}.}
\end{equation}
これは各粒子に二つの並進自由度があり、平均運動エネルギーが
$2(k_BT/2)=k_BT$ であることとも一致する。

\end{document}
